\begin{center}
{\Huge \scshape Qiong Liu, Ph.D.} \\ \vspace{4pt}
\small
\faPhone\ 475-280-1364 \quad
\href{mailto:qiongliu.work@gmail.com}{qiongliu.work@gmail.com} \quad
\href{https://www.linkedin.com/in/qiong-l-529b23196}{LinkedIn} \quad
\href{https://scholar.google.com/citations?user=SvnpHCkAAAAJ&hl=en}{Google Scholar}

Authorized to work in the U.S. without sponsorship
\end{center}

\section{Professional Summary}
\textbf{Machine Learning Researcher} with experience developing learning-based methods for spatiotemporal modeling, representation learning, signal decomposition, and robust inference from noisy real-world data. Strong background in deep learning, self-supervised learning, and data-efficient modeling, with first-author publications and patent-pending work. Experienced in taking research ideas from problem formulation through empirical validation and product-oriented collaboration.

\section{Core Expertise}
\textbf{Machine Learning:} Deep learning, self-supervised learning, multimodal learning, transformers, diffusion models, representation learning \\
\textbf{Research Areas:} Spatiotemporal modeling, temporal representation learning, signal decomposition, optimization, uncertainty modeling \\
\textbf{Programming:} Python, PyTorch, NumPy, SciPy, MATLAB \\
\textbf{Strengths:} Algorithm design, research experimentation, evaluation, cross-functional collaboration

\section{Experience}
\textbf{Canon Medical Research USA} \hfill Apr 2025 -- Present \\
\textit{AI Scientist (Reconstruction), Vernon Hills, IL}
\begin{itemize}[leftmargin=*]
\item Led development of learning-based methods for \textbf{spatiotemporal modeling} under noisy and incomplete data conditions
\item Designed \textbf{multi-signal decomposition} frameworks to separate overlapping temporal patterns for downstream modeling
\item Built end-to-end ML pipelines integrating data processing, model training, inference, and evaluation
\item Collaborated across research and product teams to translate algorithmic ideas into deployable systems
\end{itemize}

\textbf{Canon Medical Research USA} \hfill Jun 2023 -- May 2024 \\
\textit{Research Scientist Intern, Vernon Hills, IL}
\begin{itemize}[leftmargin=*]
\item Developed data-efficient deep learning models for low-signal regimes, improving robustness while preserving critical information
\item Designed evaluation frameworks to assess generalization, consistency, and quantitative reliability across datasets
\item Delivered technical findings to guide algorithm design and applied research directions
\end{itemize}

\textbf{United Imaging Healthcare} \hfill Jul 2019 -- Aug 2019 \\
\textit{Student Intern, Wuhan, China}
\begin{itemize}[leftmargin=*]
\item Supported development and validation of engineering systems for clinical applications
\end{itemize}

\section{Patents}
\textbf{Edge-Guided Deformation Field Fusion for Motion Correction} \hfill 2026 \\
Patent Pending, \textit{Lead Inventor}
\begin{itemize}[leftmargin=*]
\item Proposed a spatially adaptive learning framework to improve alignment robustness in complex dynamic data
\end{itemize}

\textbf{Automatic Data-Driven Multi-Signal Decomposition} \hfill 2025 \\
Patent Pending, \textit{Lead Inventor}
\begin{itemize}[leftmargin=*]
\item Developed a representation learning approach for separating overlapping temporal signals from unlabeled data
\end{itemize}

\section{Selected Research Projects}

\textbf{Structured Learning for Spatiotemporal Alignment}
\begin{itemize}[leftmargin=*]
\item Developed learning-based methods for alignment across sequential high-dimensional observations with challenging noise and correspondence ambiguity
\item Proposed \textbf{spatially adaptive regularization} to replace globally uniform constraints, improving robustness in low-signal and ambiguous regions
\item Studied how \textbf{structured spatial guidance} can improve optimization and generalization when standard learning objectives fail under real-world variability
\item Designed pairwise and multi-frame learning strategies for scalable temporal correspondence modeling
\end{itemize}

\textbf{Self-Supervised Learning with Physics-Informed Constraints}
\begin{itemize}[leftmargin=*]
\item Developed self-supervised temporal denoising methods that preserve meaningful latent dynamics across sequential observations
\item Incorporated \textbf{structured consistency constraints} into learning objectives, showing that physics- and domain-informed priors can substantially improve robustness in noisy, limited-data settings
\item Demonstrated improved stability and downstream reliability compared with purely data-driven baselines
\end{itemize}

\textbf{Interpretable Data-Efficient Restoration and Representation Learning}
\begin{itemize}[leftmargin=*]
\item Developed CNN-based and deep image prior approaches for restoration under heterogeneous data distributions and low-signal regimes
\item Investigated the role of \textbf{inductive bias, training data composition, and optimization dynamics} in determining model behavior and output quality
\item Analyzed model training trajectories and evolving noise structure to better understand how deep networks learn signal representations in underconstrained problems
\end{itemize}

\section{Education}
\textbf{Yale University} \hfill Aug 2020 -- May 2025 \\
Ph.D. in Biomedical Engineering \\
Thesis: Improving Image Quality and Quantification Accuracy for Static and Dynamic PET

\textbf{Huazhong University of Science and Technology} \hfill Sep 2016 -- Jun 2020 \\
B.S. in Biomedical Engineering, Minor in Computer Science \\
GPA: 3.94/4.00

\section{Selected Publications}
\textbf{First Author}
\begin{itemize}[leftmargin=*]
\item Liu Q., et al., \textit{Patlak-guided self-supervised learning for dynamic PET denoising}, \textbf{IEEE Transactions on Radiation and Plasma Medical Sciences}, 2026
\item Liu Q., et al., \textit{Population-based deep image prior for dynamic PET denoising: a data-driven approach to improve parametric quantification}, \textbf{Medical Image Analysis}, 2024
\item Liu Q., et al., \textit{A personalized deep learning denoising strategy for low-count PET images}, \textbf{Physics in Medicine \& Biology}, 2022
\end{itemize}

\section{Awards}
\textbf{Society of Nuclear Medicine and Molecular Imaging (SNMMI)}
\begin{itemize}[leftmargin=*]
\item Young Investigator Award (1st Place), 2025
\item Poster Award (2nd Place), 2024
\item Young Investigator Award (2nd Place), 2023
\end{itemize}

\textbf{Mitacs}
\begin{itemize}[leftmargin=*]
\item Globalink Research Award, 2019
\end{itemize}